% !PRINTED 291
\weberpage{291}
%<<CONT
worin
\weq{4}{\alpha\delta - \beta\gamma \equiv 1 \quad (\mathrm{mod}\ p),}
wir haben nur, wenn wir unter $\sqrt{\varDelta}$ $(\mathrm{mod}\ p)$ eine ganze
Zahl verstehen, deren Quadrat nach dem Modul $p$ kongruent
mit $\varDelta$ ist:
\[
\begin{aligned}
a &\equiv \alpha\sqrt{\varDelta}, & b &\equiv \beta\sqrt{\varDelta}\\
c &\equiv \gamma\sqrt{\varDelta}, & \partial &\equiv \delta\sqrt{\varDelta}
\end{aligned}
\quad (\mathrm{mod}\ p)
\]
zu setzen. In der Form (3) können noch die Vorzeichen von
$\alpha$, $\beta$, $\gamma$, $\delta$ gleichzeitig geändert werden.

Wir fragen nach solchen Elementen $z$, die durch eine Substitution
 von der Form (3) ungeändert bleiben. Diese werden
bestimmt durch die Kongruenz
\[
z \equiv \frac{\gamma+\delta z}{\alpha+\beta z} \quad (\mathrm{mod}\ p),
\]
oder
\weq{5}{\beta z^2 + (\alpha-\delta) z - \gamma \equiv 0 \quad (\mathrm{mod}\ p),}
eine Kongruenz, die, wenn sie nicht identisch ist, höchstens zwei
inkongruente Wurzeln hat. Um diese zu erhalten, schreiben wir,
zunächst unter der Voraussetzung, daß $\beta$ nicht durch $p$ teilbar
ist, die Kongruenz (5) so:
\[
\left(\beta z + \frac{\alpha-\delta}{2}\right)^2 \equiv \left(\frac{\alpha-\delta}{2}\right)^2 + \beta\gamma \quad (\mathrm{mod}\ p),
\]
oder mit Hilfe von (4):
\weq{6}{\left(\beta z + \frac{\alpha-\delta}{2}\right)^2 \equiv \left(\frac{\alpha+\delta}{2}\right)^2 - 1 \quad (\mathrm{mod}\ p).}
Demnach sind drei verschiedene Fälle zu unterscheiden:
\[
\mathrm{I.} \qquad \left(\frac{\alpha+\delta}{2}\right)^2 - 1 \equiv 0 \quad (\mathrm{mod}\ p);
\]
dann hat die Kongruenz (5) eine Wurzel; es gibt ein und
nur ein Element $z$, das ungeändert bleibt.
\[
\mathrm{II.} \qquad \left(\frac{\alpha+\delta}{2}\right)^2 - 1 \ \text{\gesperrt{quadratischer Rest von $p$}};
\]
die Kongruenz (5) hat zwei verschiedene Wurzeln, es gibt zwei
Elemente $z$, die ungeändert bleiben.
\[
\mathrm{III.} \qquad \left(\frac{\alpha+\delta}{2}\right)^2 - 1 \ \text{\gesperrt{quadratischer Nichtrest von $p$}};
\]
die Kongruenz (5) hat gar keine Wurzel und alle Elemente $z$
\gesperrt{werden umgesetzt.}
